% ============================================================================
% APPENDIX A: CODE IMPLEMENTATION
% ============================================================================
\chapter{Code Implementation}
\label{app:code}

This appendix provides key code excerpts from the simulation implementation. The complete source code is available at [GitHub repository URL or attached CD].

\section{Project Structure}
\label{app:structure}

\begin{lstlisting}[language=bash,caption={Project directory structure}]
ChargingAlgorithm/
├── data/
│   ├── multi_ev_config.yml           # Charge-only configuration
│   ├── multi_ev_v2g_config.yml       # V2G configuration
│   ├── office_single_building_timeseries.csv
│   └── irradiance_pvgis.csv          # PVGIS solar data
├── src/
│   ├── models/
│   │   ├── building.py               # Building model
│   │   ├── electric_vehicle.py       # EV model
│   │   ├── grid.py                   # Grid model
│   │   ├── multi_ev_charging_system.py
│   │   └── multi_ev_v2g_charging_system.py
│   ├── simulation/
│   │   ├── multi_ev_simulator.py
│   │   └── multi_ev_v2g_simulator.py
│   └── utils/
│       ├── config.py                 # Configuration loader
│       ├── data_generator.py         # Data loading utilities
│       └── visualization.py          # Plotting functions
├── scripts/
│   └── run_simulation.py             # Main entry point
└── tests/
    └── [unit tests]
\end{lstlisting}

\section{Core Models}
\label{app:core_models}

\subsection{Building Model}

\begin{lstlisting}[language=Python,caption={Building class (simplified)}]
class Building:
    def __init__(self, energy_consumption_profile, panel_area, 
                 panel_efficiency, peak_solar_irradiance,
                 battery_capacity, battery_efficiency, 
                 initial_soc, dod, irradiance_profile=None):
        self.energy_consumption_profile = energy_consumption_profile
        self.panel_area = panel_area
        self.panel_efficiency = panel_efficiency
        
        # Generate solar production profile
        if irradiance_profile:
            self.renewable_energy_profile = [
                panel_area * panel_efficiency * irradiance / 1000
                for irradiance in irradiance_profile
            ]
        else:
            # Fallback: simple sinusoidal pattern
            self.renewable_energy_profile = self._generate_solar_profile()
    
    def get_net_energy_demand(self, hour):
        """Calculate net energy demand (positive = needs grid)"""
        consumption = self.energy_consumption_profile[hour % 24]
        production = self.renewable_energy_profile[hour % 24]
        return consumption - production
\end{lstlisting}

\subsection{Electric Vehicle Model}

\begin{lstlisting}[language=Python,caption={ElectricVehicle class (simplified)}]
class ElectricVehicle:
    def __init__(self, battery_capacity, initial_soc, 
                 max_charge_rate, max_discharge_rate=0,
                 energy_per_km=0.18, dod=0.9):
        self.battery_capacity = battery_capacity
        self.soc = initial_soc
        self.max_charge_rate = max_charge_rate
        self.max_discharge_rate = max_discharge_rate
        self.min_soc = 1.0 - dod
    
    def charge(self, energy_kWh):
        """Charge the EV battery"""
        max_energy = min(
            energy_kWh,
            (1.0 - self.soc) * self.battery_capacity
        )
        actual_energy = max_energy * 0.95  # 95% efficiency
        self.soc += actual_energy / self.battery_capacity
        return actual_energy
    
    def discharge(self, energy_kWh):
        """Discharge the EV battery (V2G)"""
        max_energy = min(
            energy_kWh,
            (self.soc - self.min_soc) * self.battery_capacity
        )
        actual_energy = max_energy / 0.95  # Efficiency loss
        self.soc -= max_energy / self.battery_capacity
        return actual_energy
\end{lstlisting}

\section{Optimization Algorithms}
\label{app:algorithms}

\subsection{Simple Rule-Based Algorithm}

\begin{lstlisting}[language=Python,caption={Simple charging algorithm}]
def simple_charge_multi(self, hour):
    """Simple algorithm: Use solar first, then grid"""
    total_cost = 0
    
    # Calculate building energy surplus
    consumption = self.building.energy_consumption_profile[hour]
    production = self.building.renewable_energy_profile[hour]
    building_surplus = max(0, production - consumption)
    building_grid_draw = max(0, consumption - production)
    
    # Available grid capacity for EVs
    total_capacity = self.grid_capacity_per_hour[hour]
    available_capacity = total_capacity - building_grid_draw
    
    remaining_surplus = building_surplus
    
    # Charge each EV in sequence
    for ev_config in self.evs:
        ev = ev_config['ev']
        if hour < ev_config['arrival_time'] or \
           hour >= ev_config['target_time']:
            continue
        
        # Calculate energy needed
        energy_needed = min(
            ev.max_charge_rate,
            (ev_config['desired_soc'] - ev.soc) * ev.battery_capacity
        )
        
        # Use surplus first, then grid
        energy_from_surplus = min(energy_needed, remaining_surplus)
        energy_from_grid = min(
            energy_needed - energy_from_surplus,
            available_capacity
        )
        
        total_energy = energy_from_surplus + energy_from_grid
        if total_energy > 0:
            ev.charge(total_energy)
            cost = energy_from_grid * self.grid.get_price(hour)
            total_cost += cost
            
            # Update remaining resources
            remaining_surplus -= energy_from_surplus
            available_capacity -= energy_from_grid
    
    return True, -total_cost  # Return benefit (negative cost)
\end{lstlisting}

\subsection{Reinforcement Learning Algorithm}

\begin{lstlisting}[language=Python,caption={Q-Learning algorithm (simplified)}]
def rl_charge_multi(self, hour, episodes=2000, 
                    learning_rate=0.1, discount_factor=1.0,
                    epsilon=0.15):
    """Multi-agent Q-learning for EV charging"""
    
    # Discretize state space
    soc_bins = np.arange(0.2, 1.01, 0.1)
    actions = ['charge', 'discharge', 'standby']  # V2G
    
    # Initialize Q-table
    Q = defaultdict(lambda: np.zeros(len(actions)))
    
    # Training phase
    for episode in range(episodes):
        # Reset EV states
        ev_states = [initialize_ev_state(cfg) for cfg in self.evs]
        
        # Simulate episode
        for h in range(min_hour, max_hour):
            for ev_id, ev_state in enumerate(ev_states):
                # Discretize state
                state = discretize_state(ev_state, h)
                
                # Epsilon-greedy action selection
                if np.random.rand() < epsilon:
                    action_idx = np.random.randint(len(actions))
                else:
                    action_idx = np.argmax(Q[state])
                
                # Execute action and calculate reward
                reward = execute_action_and_get_reward(
                    ev_state, actions[action_idx], h
                )
                
                # Update Q-value
                next_state = get_next_state(ev_state, h)
                best_next_q = np.max(Q[next_state])
                Q[state][action_idx] += learning_rate * (
                    reward + discount_factor * best_next_q 
                    - Q[state][action_idx]
                )
    
    # Execution phase: Use learned policy
    for ev_config in self.evs:
        state = get_current_state(ev_config, hour)
        action_idx = np.argmax(Q[state])
        execute_action(ev_config, actions[action_idx])
    
    return True, total_benefit
\end{lstlisting}

\section{Simulation Orchestration}
\label{app:simulation}

\begin{lstlisting}[language=Python,caption={Main simulation loop}]
def run_multi_ev_v2g_simulation(config):
    """Run complete multi-EV V2G simulation"""
    
    # Initialize components
    building = Building(...)
    grid = Grid(...)
    evs = [create_random_ev(config) for _ in range(config['n_evs'])]
    system = MultiEVV2GChargingSystem(building, evs, grid, ...)
    
    # Define methods to compare
    methods = {
        'simple': system.simple_charge_multi,
        'rl': lambda h: system.rl_charge_multi(h, episodes=15000),
    }
    
    # Run simulation
    results = {name: defaultdict(list) for name in methods}
    
    for hour in range(24):
        for name, func in methods.items():
            # Reset EVs to same initial state
            reset_ev_states(evs, results[name], hour)
            
            # Execute algorithm
            acted, benefit = func(hour)
            
            # Record results
            results[name]['benefit'].append(benefit)
            results[name]['ev_soc'].append([ev['ev'].soc for ev in evs])
    
    # Visualize and analyze
    visualize_results(results, config)
    
    return results
\end{lstlisting}

\section{Configuration Management}
\label{app:config}

\subsection{Configuration File Format}

\begin{lstlisting}[caption={Example configuration file (multi\_ev\_v2g\_config.yml)}]
# Building Configuration
panel_area: 150
panel_efficiency: 0.20
building_battery_capacity: 80
building_battery_efficiency: 0.95

# EV Fleet Configuration
n_evs: 15
ev_battery_capacity: 100
ev_initial_soc_range: [0.2, 0.3]
ev_desired_soc_range: [0.7, 0.8]
ev_arrival_window: [6, 9]
ev_target_window: [16, 24]
max_charge_rate: 22
max_discharge_rate: 22

# Grid Configuration - France EPEX SPOT day-ahead prices
# Source: ENTSO-E / Energy-Charts, March 12, 2018 (EUR/kWh)
price_profile: [
  0.0285, 0.0198, 0.0120, 0.0096,  # 00-03h
  0.0190, 0.0420, 0.0447, 0.0483,  # 04-07h
  0.0481, 0.0460, 0.0444, 0.0450,  # 08-11h
  0.0415, 0.0384, 0.0372, 0.0376,  # 12-15h
  0.0422, 0.0502, 0.0620, 0.0552,  # 16-19h
  0.0432, 0.0412, 0.0368, 0.0350   # 20-23h
]
# V2G sell-back at 80% of purchase price
sell_price_profile: [
  0.0228, 0.0158, 0.0096, 0.0076,  # 00-03h
  0.0152, 0.0336, 0.0357, 0.0387,  # 04-07h
  0.0384, 0.0368, 0.0355, 0.0360,  # 08-11h
  0.0332, 0.0307, 0.0298, 0.0301,  # 12-15h
  0.0338, 0.0402, 0.0496, 0.0442,  # 16-19h
  0.0346, 0.0330, 0.0294, 0.0280   # 20-23h
]
grid_capacity_per_hour: [80, 80, 80, ...]

# Algorithm Parameters
rl_episodes: 15000
rl_lr: 0.15
rl_gamma: 1.0
rl_epsilon: 0.10
\end{lstlisting}

\section{Data Loading Utilities}
\label{app:data_loading}

\begin{lstlisting}[language=Python,caption={PVGIS data loader}]
def load_irradiance_pvgis(filepath, date_str):
    """
    Load PVGIS irradiance data for a specific date.
    
    Args:
        filepath: Path to PVGIS CSV file
        date_str: Date in format 'YYYYMMDD'
    
    Returns:
        List of 24 hourly irradiance values (W/m²)
    """
    import pandas as pd
    
    df = pd.read_csv(filepath)
    
    # Filter for specific date
    df['date'] = pd.to_datetime(df['time']).dt.strftime('%Y%m%d')
    df_day = df[df['date'] == date_str]
    
    # Extract hourly irradiance
    irradiance = df_day['G(i)'].values  # Global irradiance on inclined plane
    
    # Ensure 24 hours
    if len(irradiance) < 24:
        irradiance = np.pad(irradiance, (0, 24 - len(irradiance)))
    
    return irradiance[:24].tolist()
\end{lstlisting}

\section{Visualization Functions}
\label{app:visualization}

\begin{lstlisting}[language=Python,caption={Results visualization (simplified)}]
def visualize_multi_ev_v2g(results, config, evs, system):
    """Create comprehensive visualization of simulation results"""
    
    fig, axes = plt.subplots(4, 1, figsize=(16, 12))
    
    # Plot 1: EV SoC evolution
    for method in ['simple', 'rl']:
        socs = np.array(results[method]['ev_soc'])
        avg_soc = socs.mean(axis=1)
        axes[0].plot(results['hours'], avg_soc, 
                    label=method.upper(), linewidth=2)
    
    axes[0].set_ylabel('Average SoC')
    axes[0].set_title('EV Fleet State of Charge')
    axes[0].legend()
    axes[0].grid(True)
    
    # Plot 2: Grid usage
    for method in ['simple', 'rl']:
        axes[1].plot(results['hours'], 
                    results[method]['grid_usage'],
                    label=method.upper(), marker='o')
    
    # Plot grid capacity limit
    capacity = [config['grid_capacity_per_hour'][h] 
                for h in results['hours']]
    axes[1].plot(results['hours'], capacity, 
                'k--', label='Grid Capacity')
    
    axes[1].set_ylabel('Grid Power (kW)')
    axes[1].set_title('Grid Power Consumption')
    axes[1].legend()
    axes[1].grid(True)
    
    # Plot 3: Hourly benefit
    # [Similar plotting code...]
    
    # Plot 4: Cumulative benefit
    # [Similar plotting code...]
    
    plt.tight_layout()
    plt.savefig('simulation_results_multi_v2g.png', dpi=300)
\end{lstlisting}

\section{Usage Instructions}
\label{app:usage}

\subsection{Running Simulations}

\begin{lstlisting}[language=bash,caption={Command line usage}]
# Install dependencies
pip install -r requirements.txt

# Run charge-only simulation
python scripts/run_simulation.py  # Set SIMULATION_TYPE = "multi_ev"

# Run V2G simulation
python scripts/run_simulation.py  # Set SIMULATION_TYPE = "multi_ev_v2g"

# Run tests
pytest tests/
\end{lstlisting}

\subsection{Modifying Parameters}

To experiment with different scenarios:

\begin{enumerate}
    \item Edit configuration file: \texttt{data/multi\_ev\_v2g\_config.yml}
    \item Modify parameters (fleet size, prices, capacities, etc.)
    \item Run simulation: \texttt{python scripts/run\_simulation.py}
    \item Results saved to: \texttt{scripts/simulation\_results\_*.png}
\end{enumerate}

\section{Dependencies}
\label{app:dependencies}

\begin{lstlisting}[caption={requirements.txt}]
numpy>=1.21.0
matplotlib>=3.4.0
pyyaml>=5.4.0
pandas>=1.3.0
pytest>=7.0.0            # For testing
\end{lstlisting}

\section{Testing}
\label{app:testing}

\begin{lstlisting}[language=Python,caption={Example unit test}]
import pytest
from src.models.electric_vehicle import ElectricVehicle

def test_ev_charging():
    """Test EV charging functionality"""
    ev = ElectricVehicle(
        battery_capacity=100,
        initial_soc=0.3,
        max_charge_rate=22
    )
    
    # Charge for 1 hour at max rate
    energy = ev.charge(22)
    
    # Check energy charged (with efficiency)
    assert abs(energy - 20.9) < 0.1  # 22 * 0.95
    
    # Check SoC increase
    expected_soc = 0.3 + 20.9 / 100
    assert abs(ev.soc - expected_soc) < 0.01

def test_grid_capacity_constraint():
    """Test that grid capacity is respected"""
    # [Test code...]
    pass

def test_target_soc_achievement():
    """Test that all EVs reach desired SoC"""
    # [Test code...]
    pass
\end{lstlisting}

\section{Performance Profiling}
\label{app:profiling}

\begin{lstlisting}[language=Python,caption={Algorithm timing comparison}]
import time

def benchmark_algorithms():
    """Compare execution time of different algorithms"""
    
    algorithms = {
        'simple': system.simple_charge_multi,
        'rl': lambda h: system.rl_charge_multi(h, episodes=2000),
        'mpc': lambda h: system.mpc_charge(h, horizon=6),
        'pso': lambda h: system.pso_charge(h, n_particles=30),
        'milp': system.milp_charge,
    }
    
    results = {}
    
    for name, func in algorithms.items():
        start = time.time()
        
        for hour in range(24):
            func(hour)
        
        elapsed = time.time() - start
        results[name] = elapsed
        print(f"{name}: {elapsed:.2f} seconds")
    
    return results
\end{lstlisting}

\section{Code Availability}
\label{app:availability}

The complete source code for this thesis is available at:

\begin{itemize}
    \item \textbf{Repository:} [GitHub URL]
    \item \textbf{License:} [Specify license, e.g., MIT, GPL]
    \item \textbf{Documentation:} See \texttt{README.md} for detailed usage instructions
    \item \textbf{Contact:} [Your email] for questions or collaboration
\end{itemize}

\section{Reproducibility}
\label{app:reproducibility}

To reproduce the results in this thesis:

\begin{enumerate}
    \item Clone the repository
    \item Install dependencies: \texttt{pip install -r requirements.txt}
    \item Download data files (included in repository)
    \item Run: \texttt{python scripts/run\_simulation.py}
    \item Results will be generated in \texttt{scripts/} folder
\end{enumerate}

\textbf{Note:} RL results may vary slightly due to random initialization. Use \texttt{np.random.seed(42)} for deterministic results.
